% Part: proof-theory
% Chapter: cut-elimination
% Section: interpolation

\documentclass[../../../include/open-logic-section]{subfiles}

\begin{document}

\olfileid{pt}{cut}{itp}

\olsection{د مايېهارا مرستندويه قضيه او د کرېګ د منځګړيتوب قضيه}

د منځګړيتوب قضيه، چې ويليم کرېګ وړاندې کړې ده، وايي: که $!A
\Entails !B$ وي، نو داسې !!a{sentence}~$!C$ (\emph{منځګړی}) شته
چې $!A \Entails !C$ او $!C \Entails !B$ وي۔ منځګړی~$!C$ داسې
غوره کېدے شي چې په~$!C$ کښې راغلي ټول !!{constant}s او !!{predicate}s
په~$!A$ او~$!B$ دواړو کښې راغلي وي۔ (فرض کوو چې هېڅ !!{function}s
نۀ کارېږي۔) له دې قيد پرته، $!A$~او $!B$ په اسانۍ منځګړي کېدے شي۔

د منځګړيتوب قضيه د کلاسيکي لومړۍ درجې منطق دپاره صدق کوي او
دا د لومړۍ درجې منطق «وروستے مهم کشف شوے خاصيت» بلل شوے دے۔
په مدل پوهنه او اتومات استدلال کښې مهم استعمالونه لري۔ د دې
لومړنے محرک او استعمال د علم په فلسفه کښې ؤ، ځکه د دې په وسيله
کتلے شو چې يوه تيوري په کومو لومړنيو مفاهيمو بديهياتي کېدے شي
يا نۀ شي کېدے۔ دا د بېت د تعريف‌وړتيا قضيه هم لازموي: که يوه
تيوري يو مفهوم په ضمني توګه تعريفولے شي، نو په څرګند ډول ئې هم
تعريفولے شي۔

لکه په رياضيکي منطق کښې ډېرې پايلې، د منځګړيتوب قضيه هم د
مدل پوهنې او هم د ثبوت پوهنې په طريقو ثابتېدے شي۔ د ثبوت پوهنې
د طريقو ګټه دا ده چې د منځګړي شتون ښيي او داسې الګوريتم هم راکوي
چې منځګړی د دې خبرې له !!a{proof} څخه محاسبه کوي چې $!A \Entails !B$
دے؛ مثلاً په~\Log{G3c} کښې د $!A \Sequent !B$ له !!a{proof} څخه۔

دې $\Lang L(!A)$ په~$!A$ کښې د !!{constant}s او !!{predicate}s
سټ وي او $\Lang L(\Gamma) = \bigcup \Setabs{\Lang L(!A)}{!A \in
\Gamma}$ دې وي۔ په دې ټوله برخه کښې فرض کوو چې له داسې
!!{formula}s سره کار کوو چې !!{function}s نۀ لري۔

\begin{thm}[د کرېګ د منځګړيتوب قضيه]\ollabel{thm:interpolation}
دې د~$!A$ !!{language} $\Lang L(!A)$ او د $!B$ ژبه
$\Lang L(!B)$ وي۔ که $\Log{G3c} \Proves !A \Sequent !B$ وي،
نو په !!{language}~$\Lang L(!C) \subseteq \Lang
L(!A) \cap \Lang L(!B)$ کښې داسې !!a{formula}~$!C$ شته چې
$\Log{G3c} \Proves !A \Sequent !C$ او
$\Log{G3c} \Proves !C \Sequent !B$ وي۔
\end{thm}

د منځګړيتوب د قضيې يوه عمومي بڼه ثابتوو، څو په~\Log{G3c} کښې
د !!{proof}s له جوړښت څخه کار واخلو۔ د دې دپاره د سېکوېنټونو
مقدم او تالي هر يو په دوو برخو وېشل شوي وګڼئ۔
$\maeh{\Gamma_1}{\Gamma_2} \Sequent
\maeh{\Delta_1}{\Delta_2}$ داسې ليکو چې $\Gamma_1$ او $\Delta_1$
لومړۍ برخې ته او $\Gamma_2$ او $\Delta_2$ دويمې برخې ته اړه ولري۔
بيا د منځګړيتوب قضيه له لاندې مرستندويې قضيې څخه نېغ په نېغه راځي۔

\begin{lem}[د مايېهارا مرستندويه قضيه]\ollabel{lem:maehara}
که $\Log{G3c}
\Proves \maeh{\Gamma_1}{\Gamma_2} \Sequent \maeh{\Delta_1}{\Delta_2}$
وي، نو په !!{language}~$\Lang L(!C)
\subseteq \Lang L(\Gamma_1,
\Delta_1) \cap \Lang L(\Gamma_2, \Delta_2)$ کښې داسې !!a{formula}~$!C$
شته چې $\Log{G3c}
\Proves \Gamma_1 \Sequent \Delta_1, !C$ او $\Log{G3c} \Proves !C,
\Gamma_2 \Sequent \Delta_2$ وي۔
\end{lem}

\begin{proof}
  فرض کړئ چې $\pi \Proves \maeh{\Gamma_1}{\Gamma_2} \Sequent
  \maeh{\Delta_1}{\Delta_2}$ دے۔ دعوه په
  $n=\pheight{\pi}$ باندې په استقرا ثابتوو۔

  که $n=0$ وي، نو $\maeh{\Gamma_1}{\Gamma_2} \Sequent
  \maeh{\Delta_1}{\Delta_2}$ يو بديهي اصل دے او کوم اټومي
  !!{formula}~$!A$ په $\Gamma_1, \Gamma_2$ او $\Delta_1,
  \Delta_2$ دواړو کښې راغلے دے۔ څلور حالتونه دي: په کيڼ اړخ
  کښې $!A$~لومړۍ که دويمې برخې ته اړه لري، او په ښي اړخ کښې
  لومړۍ که دويمې برخې ته۔
  \begin{enumerate}
    \item $!A \in \Gamma_1$ او $!A \in \Delta_1$۔ داسې~$!C$
    غواړو چې $\Gamma_1 \Sequent \Delta_1, !C$ او $!C, \Gamma_2
    \Sequent \Delta_2$ دواړه !!{provable} وي۔ لومړنے له وړاندې
    بديهي اصل دے، که~$!C$ هر ډول غوره کړو، ځکه $!A$ په کيڼ او ښي
    دواړو اړخونو کښې شته۔ د دې دپاره چې $!C, \Gamma_2
    \Sequent \Delta_2$ خامخا !!{provable} وي، دلته عمومي خوندي
    انتخاب دا دے چې $!C \ident \lfalse$ واخلو۔
    \item $!A \in \Gamma_1$ او $!A \in \Delta_2$۔ نو $\Gamma_1
    \Sequent \Delta_1, !A$ او $!A, \Gamma_2 \Sequent \Delta_2$
    دواړه بديهي اصلونه دي؛ ځکه $!C \ident !A$ اخستلے شو۔
    \item $!A \in \Gamma_2$ او $!A \in \Delta_1$۔ نو $!A, \Gamma_1
    \Sequent \Delta_1$ او $\Gamma_2 \Sequent \Delta_2, !A$ به
    بديهي اصلونه وي، خو~$!A$ ئې په ناسم اړخونو کښې دے۔ نو $!A$
    منځګړی نۀ دے۔ خو د $\RightR{\lnot}$ او $\LeftR{\lnot}$ په
    وسيله $\Gamma_1
    \Sequent \Delta_1, \lnot!A$ او $\lnot !A, \Gamma_2 \Sequent
    \Delta_2$ !!{prove} کولے شو۔
    \item $!A \in \Gamma_2$ او $!A \in \Delta_2$۔ تمرين۔
  \end{enumerate}
  هغه حالتونه هم تمرين ته پرېږدو چې $\maeh{\Gamma_1}{\Gamma_2} \Sequent
  \maeh{\Delta_1}{\Delta_2}$ ځکه بديهي اصل وي چې $\lfalse \in
  \Gamma_1$ يا $\lfalse \in \Gamma_2$ دے۔

  که $n>0$ وي، $\pi$ په يوه منطقي استنتاج ختمېږي۔ د کارول شوې
  قاعدې او د اصلي فارمول د برخې له مخې بېل حالتونه ګورو۔ په هر
  حالت کښې دا استقرايي فرض کارولے شو چې مرستندويه قضيه د ټولو
  هغو !!{proof}s~$\pi_1$ دپاره صدق کوي چې
  $\pheight{\pi_1} < n$ وي، که د~$\pi_1$ وروستے سېکوېنټ هر ډول
  په دوو برخو وېشل شوے وي۔ څو ځانګړي حالتونه وګورو۔
  \begin{enumerate}
    \item $\pi$ په $\RightR{\lnot}$ ختمېږي او اصلي فارمول ئې~$\lnot
    !A$ دے۔ دوه حالتونه دي:~$\lnot !A$ لومړۍ که دويمې برخې ته
    اړه لري؛ د دې له مخې $\pi$ په دې ختمېږي:
\[
    \AxiomC{}
    \RightLabel{$\pi_1$}
    \Deduce$\maeh{!A, \Gamma_1}{\Gamma_2} \fCenter
      \maeh{\Delta_1}{\Delta_2}$
    \RightLabel{\RightR{\lnot}}
    \UnaryInf$\maeh{\Gamma_1}{\Gamma_2} \fCenter
      \maeh{\Delta_1, \lnot !A}{\Delta_2}$
    \DisplayProof
    \]
    يا
\[
    \AxiomC{}
    \RightLabel{$\pi_1$}
    \Deduce$\maeh{\Gamma_1}{!A, \Gamma_2} \fCenter
      \maeh{\Delta_1}{\Delta_2}$
    \RightLabel{\RightR{\lnot}}
    \UnaryInf$\maeh{\Gamma_1}{\Gamma_2} \fCenter
      \maeh{\Delta_1}{\Delta_2, \lnot !A}$
    \DisplayProof
    \]
    پام وکړئ: که اصلي !!{formula}~$\lnot !A$ لومړۍ برخې ته اړه
    لري، نو په مقدمه کښې مو فعال !!{formula}~$!A$ هم لومړۍ برخې
    ته ورکړے دے۔ دا انتخاب زمونږ په لاس کښې دے۔ استقرايي فرض
    به هله هم کارېدۀ چې فعال فارمول مو بلې برخې ته ورکړے وے، خو
    بيا به مو منځګړی نۀ شو موندلے (تمرين: وې ازمويئ)۔

    لومړے هغه حالت وګورئ چې~$\lnot !A$ لومړۍ برخې ته اړه لري۔
    د استقرايي فرض له مخې، د~$\pi_1$ د وروستي سېکوېنټ منځګړی
    شته؛ يعنې داسې !!a{formula}~$!C_1$ چې دواړه لاندې سېکوېنټونه
\begin{align*}
    !A, \Gamma_1 & \Sequent \Delta_1, !C_1 &&\text{او} &
    !C_1, \Gamma_2 & \Sequent \Delta_2
    \intertext{د ثابتولو وړ دي۔ مونږ داسې !!a{formula}~$!C$ غواړو چې
    لاندې دواړه سېکوېنټونه !!{provable} وي:}
    \Gamma_1 & \Sequent \Delta_1, \lnot !A, !C &&\text{او} &
    !C, \Gamma_2 & \Sequent \Delta_2
    \end{align*}
    څرګنده ده چې $!C \ident !C_1$ اخستلے شو: استقرايي فرض له
    وړاندې وايي چې د ښي اړخ سېکوېنټ !!{provable} دے؛ د کيڼ اړخ
    سېکوېنټ له هغه سېکوېنټ څخه په~$\RightR{\lnot}$ راځي چې له
    استقرايي فرضه مو تر لاسه کړے ؤ۔ استقرايي فرض دا هم وايي چې
    $\Lang L(!C_1) \subseteq \Lang L(!A, \Gamma_1,
    \Delta_1) \cap \Lang L(\Gamma_2, \Delta_2)$۔ ځکه چې $\Lang L(\lnot
    !A) = \Lang L(!A)$ او $\Lang L(!C) = \Lang L(!C_1)$ دي، نو
    $\Lang L(!C_1) \subseteq \Lang L(\Gamma_1, \Delta_1, \lnot !A) \cap \Lang
    L(\Gamma_2, \Delta_2)$ لرو۔

    په دويم حالت کښې وضع لږه بېله ده، ځکه~$\lnot !A$ دويمې
    برخې ته اړه لري۔ په مقدمه کښې فعال !!{formula} هم دويمې
    برخې ته ورکوو او استقرايي فرض کاروو، څو لاندې سېکوېنټونه
\begin{align*}
    \Gamma_1 & \Sequent \Delta_1, !C_1 &&\text{او} &
    !C_1, !A, \Gamma_2 & \Sequent \Delta_2
    \intertext{د ثابتولو وړ وي۔ که بيا د دويم سېکوېنټ دپاره د \RightR{\lnot} قاعده وکاروو، د لاندې سېکوېنټونو !!{proof}s تر لاسه کوو}
    \Gamma_1 & \Sequent \Delta_1, !C_1 &&\text{او} &
    !C_1, \Gamma_2 & \Sequent \Delta_2, \lnot !A
    \end{align*}
    دا بيا عين هغه سېکوېنټونه دي چې د $!C_1$ د منځګړي کېدو
    دپاره ئې !!{provable} کېدل غواړو؛ بيا $!C \ident !C_1$ اخلو۔
    لکه مخکې، $\Lang
    L(!C_1) \subseteq \Lang L(\Gamma_1, \Delta_1) \cap \Lang L(\Gamma_2,
    \Delta_2, \lnot !A)$ دے۔
    \item $\pi$ په $\RightR{\lif}$ ختمېږي او اصلي !!{formula}
    ئې~$!A \lif !B$ دے۔ بيا دوه حالتونه دي: $!A \lif
    !B$ لومړۍ که دويمې برخې ته اړه لري۔ !!{proof}~$\pi$ په دې
    ختمېږي:
\[
    \AxiomC{}
    \RightLabel{$\pi_1$}
    \Deduce$\maeh{!A, \Gamma_1}{\Gamma_2} \fCenter
      \maeh{\Delta_1, !B}{\Delta_2}$
    \RightLabel{\RightR{\lif}}
    \UnaryInf$\maeh{\Gamma_1}{\Gamma_2} \fCenter
      \maeh{\Delta_1, !A \lif !B}{\Delta_2}$
    \DisplayProof
    \]
    يا
\[
    \AxiomC{}
    \RightLabel{$\pi_1$}
    \Deduce$\maeh{\Gamma_1}{!A, \Gamma_2} \fCenter
      \maeh{\Delta_1}{\Delta_2, !B}$
    \RightLabel{\RightR{\lif}}
    \UnaryInf$\maeh{\Gamma_1}{\Gamma_2} \fCenter
      \maeh{\Delta_1}{\Delta_2, !A \lif !B}$
    \DisplayProof
    \]
    بيا مو په مقدمه کښې فعال !!{formula}s هماغې برخې ته ورکړي
    دي چې په پايله کښې اصلي !!{formula} ورته اړه لري۔ په لومړي
    حالت کښې استقرايي فرض کاروو او لاندې !!{proof}s تر لاسه کوو:
\begin{align*}
    !A, \Gamma_1 & \Sequent \Delta_1, !B, !C_1 &&\text{او} & !C_1, \Gamma_2 & \Sequent \Delta_2
    \intertext{د کيڼ اړخ سېکوېنټ د \RightR{\lif} دپاره مناسبه مقدمه ده۔ نو لاندې سېکوېنټونه !!{provable} دي:}
    \Gamma_1 & \Sequent \Delta_1, !A \lif !B, !C_1 &&\text{او} & !C_1, \Gamma_2 & \Sequent \Delta_2
    \end{align*}
    نو $!C_1$ د~$\pi$ د وروستي سېکوېنټ دپاره هم منځګړی دے او
    يو ځل بيا $!C \ident !C_1$ اخستلے شو۔

    بل حالت هم په همدې ډول څېړل کېږي۔
    \item $\pi$ په $\LeftR{\lif}$ ختمېږي او اصلي !!{formula}
    ئې~$!A \lif !B$ دے۔ که $!A \lif !B$ لومړۍ برخې ته اړه
    ولري، نو !!{proof}~$\pi$ په دې ختمېږي:
\[
    \AxiomC{}
    \RightLabel{$\pi_1$}
    \Deduce$\maeh{\Gamma_1}{\Gamma_2} \fCenter
      \maeh{\Delta_1, !A}{\Delta_2}$
    \AxiomC{}
    \RightLabel{$\pi_2$}
    \Deduce$\maeh{!B, \Gamma_1}{\Gamma_2} \fCenter
      \maeh{\Delta_1}{\Delta_2}$
    \BinaryInf$\maeh{!A \lif !B, \Gamma_1}{\Gamma_2} \fCenter
      \maeh{\Delta_1}{\Delta_2}$
    \DisplayProof
    \]
    استقرايي فرض اوس څلور !!{provable} سېکوېنټونه راکوي؛ دوه
    د کيڼې مقدمې د منځګړي $!C_1$ دپاره او دوه د دويمې مقدمې
    د منځګړي $!C_2$ دپاره:
\begin{align*}
    \Gamma_1 & \Sequent \Delta_1, !A, !C_1 &&\text{(1)} &
    !C_1, \Gamma_2 & \Sequent \Delta_2 && \text{(2)}\\
    !B, \Gamma_1 & \Sequent \Delta_1, !C_2 &&\text{(3)} &
    !C_2, \Gamma_2 & \Sequent \Delta_2 && \text{(4)}
    \end{align*}
    سېکوېنټونه (1) او (3) د~\LeftR{\lif} مقدمې کېدے شي، که
    کمزوري کول وکاروو او د~(1) په ښي اړخ $!C_2$ او د~(3) په
    ښي اړخ $!C_1$ ورزيات کړو:
\[
    \AxiomC{}
    \Deduce$\Gamma_1 \fCenter \Delta_1, !A, !C_1, !C_2$
    \AxiomC{}
    \Deduce$!B, \Gamma_1 \fCenter \Delta_1, !C_1, !C_2$
    \RightLabel{\LeftR{\lif}}
    \BinaryInf$!A \lif !B, \Gamma_1 \fCenter \Delta_1, !C_1, !C_2$
    \DisplayProof
    \]
    د $\RightR{\lor}$ په کارولو دا سېکوېنټ تر لاسه کوو:
\[
    !A \lif !B, \Gamma_1 \fCenter \Delta_1, !C_1 \lor !C_2
    \]
    دا ښيي چې $!C_1 \lor !C_2$ په دې حالت کښې زمونږ منځګړی
    کېدے شي۔ او په رښتيا، بل اړين سېکوېنټ له پاتې سېکوېنټونو
    (2) او~(4) څخه !!{prove} کولے شو:
\[
    \AxiomC{}
    \Deduce$!C_1, \Gamma_2 \fCenter \Delta_2$
    \AxiomC{}
    \Deduce$!C_2, \Gamma_2 \fCenter \Delta_2$
    \RightLabel{\LeftR{\lor}}
    \BinaryInf$!C_1 \lor !C_2, \Gamma_2 \fCenter \Delta_2$
    \DisplayProof
    \]
    دا کتنه پاتې ده چې $!C_1 \lor !C_2$ په سمه ژبه کښې دے۔
    استقرايي فرض راکوي چې
\begin{align*}
    \Lang L(!C_1) & \subseteq
    \Lang L(\Gamma_1, \Delta_1, !A) \cap \Lang L(\Gamma_2, \Delta_2) &\text{او}\\
    \Lang L(!C_2) & \subseteq
    \Lang L(!B, \Gamma_1, \Delta_1) \cap \Lang L(\Gamma_2, \Delta_2)
    \intertext{ځکه چې}
    \Lang L(!A) & \subseteq \Lang L(!A \lif !B) &\text{او}\\
    \Lang L(!B) & \subseteq \Lang L(!A \lif !B)
    \intertext{دواړه لاندې اړيکې لرو}
    \Lang L(!C_1) & \subseteq
    \Lang L(\Gamma_1, \Delta_1, !A \lif !B) \cap \Lang L(\Gamma_2, \Delta_2) &\text{او}\\
    \Lang L(!C_2) & \subseteq
    \Lang L(\Gamma_1, \Delta_1, !A \lif !B) \cap \Lang L(\Gamma_2, \Delta_2)
    \intertext{او په پايله کښې، $\Lang L (!C_1 \lor !C_2) = {}$}
    \Lang L(!C_1) \cup \Lang L(!C_2) & \subseteq
    \Lang L(\Gamma_1, \Delta_1, !A \lif !B) \cap \Lang L(\Gamma_2, \Delta_2),
    \end{align*}
    لکه چې غوښتل شوي وو۔

    که $!A \lif !B$ دويمې برخې ته اړه لري، وضع بېله ده، خو
    په ورته ډول څېړل کېږي۔ استقرايي فرض بيا څلور !!{provable}
    سېکوېنټونه راکوي:
\begin{align*}
    \Gamma_1 & \Sequent \Delta_1, !C_1 &&\text{(1)} &
    !C_1, \Gamma_2 & \Sequent \Delta_2, !A && \text{(2)}\\
    \Gamma_1 & \Sequent \Delta_1, !C_2 &&\text{(3)} &
    !C_2, !B, \Gamma_2 & \Sequent \Delta_2 && \text{(4)}
    \end{align*}
    اوس سېکوېنټونه (2) او (4)، له ځينو د کمزوري کولو په وسيله
    زياتو شوو !!{formula}s سره، د~\LeftR{\lif} مقدمې جوړوي:
\[
    \AxiomC{}
    \Deduce$!C_1, !C_2, \Gamma_2 \fCenter \Delta_2, !A$
    \AxiomC{}
    \Deduce$!B, !C_1, !C_2, \Gamma_2 \fCenter \Delta_2$
    \RightLabel{\LeftR{\lif}}
    \BinaryInf$!A \lif !B, !C_1, !C_2, \Gamma_2 \fCenter \Delta_2$
    \DisplayProof
    \]
    بيا غواړو دا داسې سېکوېنټ کړو چې له $!C_1$ او $!C_2$ څخه
    جوړ يو واحد !!{formula} ولري؛ خو دا ځل ئې په مقدم کښې
    غواړو، ځکه اوس له دويمې برخې سره کار کوو۔ د $\LeftR{\land}$
    په کارولو دا ترلاسه کېږي؛ منځګړی بايد $!C \ident
    (!C_1 \land !C_2)$ واخلو۔ پاتې سېکوېنټونه (1) او (3) هم
    د بلې برخې اړوند سېکوېنټ د !!{prove} کولو دپاره کارولے شو:
\[
    \AxiomC{}
    \Deduce$\Gamma_1 \fCenter \Delta_1, !C_1$
    \AxiomC{}
    \Deduce$\Gamma_1 \fCenter \Delta_1, !C_2$
    \RightLabel{\RightR{\land}}
    \BinaryInf$\Gamma_1 \fCenter \Delta_1, !C_1 \land !C_2$
    \DisplayProof
    \]
    لکه مخکې، $\Lang L(!C_1 \land !C_2) \subseteq \Lang
    L(\Gamma_1, \Delta_1) \cap \Lang L(\Gamma_2, \Delta_2, !A \lif
    !B)$ لرو۔

    \item $\pi$ په $\RightR{\lforall}$ ختمېږي او اصلي !!{formula}
    ئې~$\lforall[x][!A(x)]$ دے:
\[
    \AxiomC{}
    \RightLabel{$\pi_1$}
    \Deduce$\maeh{\Gamma_1}{\Gamma_2} \fCenter
    \maeh{\Delta_1, !A(c)}{\Delta_2}$
    \RightLabel{\RightR{\lforall}}
    \UnaryInf$\maeh{\Gamma_1}{\Gamma_2} \fCenter
      \maeh{\Delta_1, \lforall[x][!A(x)]}{\Delta_2}$
    \DisplayProof
    \]
    يا
\[
    \AxiomC{}
    \RightLabel{$\pi_1$}
    \Deduce$\maeh{\Gamma_1}{\Gamma_2} \fCenter
      \maeh{\Delta_1}{\Delta_2, !A(c)}$
    \RightLabel{\RightR{\lforall}}
    \UnaryInf$\maeh{\Gamma_1}{\Gamma_2} \fCenter
      \maeh{\Delta_1}{\Delta_2, \lforall[x][!A(x)]}$
    \DisplayProof
    \]
    په لومړي حالت کښې استقرايي فرض د $\Gamma_1 \Sequent \Delta_1, !A(c), !C_1$
    او $!C_1, \Gamma_2
    \Sequent \Delta_2$ !!{proof}s راکوي۔ په لومړي سېکوېنټ
    $\RightR{\lforall}$ کارولے شو، څو $\Gamma_1 \Sequent \Delta_1, \lforall[x][!A(x)], !C_1$
    تر لاسه کړو۔ (د ځانګړي متغير شرط پوره دے، ځکه د~$\pi$ په
    پای کښې د \RightR{\lforall} په قاعده کښې له وړاندې پوره ؤ۔)
    نو $!C_1$ به منځګړی وي، که دا شرط پوره وي چې $\Lang
    L(!C_1) \subseteq \Lang L(\Gamma_1, \Delta_1, \lforall[x][!A(x)])
    \cap \Lang L(\Gamma_2, \Delta_2)$۔ د استقرايي فرض له مخې،
    $\Lang L(!C_1) \subseteq \Lang L(\Gamma_1,
    \Delta_1, !A(c)) \cap \Lang L(\Gamma_2, \Delta_2)$ لرو۔ نو
    د~$c$ په هکله بايد پام وکړو: آيا $c$ په~$!C_1$ کښې راتلے
    شي؟ نۀ شي راتلے۔ که $c$ په $!C_1$ کښې راشي، نو هم
    $c \in \Lang L(\Gamma_1, \Delta_1,
    !A(c))$ (چې رښتيا ده) او هم $c \in \Lang L(\Gamma_2,
    \Delta_2)$ به وي۔ خو بيا به $c$ په~$\pi$ کښې د
    $\RightR{\lforall}$ د استنتاج په پايله کښې راشي او د ځانګړي
    متغير شرط به مات کړي۔

    بل حالت هم ورته دے۔

    \item $\pi$ په $\RightR{\lexists}$ ختمېږي او اصلي
    !!{formula} ئې~$\lexists[x][!A(x)]$ دے۔ بيا دوه حالتونه دي۔
    هغه حالت ګورو چې $\lexists[x][!A(x)]$ لومړۍ برخې ته اړه لري؛
    بل حالت تمرين ته پرېږدو۔

    په لومړي حالت کښې $\pi$ په دې ختمېږي:
\[
    \AxiomC{}
    \RightLabel{$\pi_1$}
    \Deduce$\maeh{\Gamma_1}{\Gamma_2} \fCenter
      \maeh{\Delta_1, \lexists[x][!A(x)], !A(t)}{\Delta_2}$
    \RightLabel{\RightR{\lexists}}
    \UnaryInf$\maeh{\Gamma_1}{\Gamma_2} \fCenter
      \maeh{\Delta_1, \lexists[x][!A(x)]}{\Delta_2}$
    \DisplayProof
    \]
    استقرايي فرض د $\Gamma_1 \Sequent
    \Delta_1, \lexists[x][!A(x)], !A(t), !C_1$ او $!C_1, \Gamma_2
    \Sequent \Delta_2$ !!{proof}s راکوي۔ په لومړي سېکوېنټ
    $\RightR{\lexists}$ کارولے شو، څو $\Gamma_1 \Sequent \Delta_1, \lexists[x][!A(x)], !C_1$
    تر لاسه کړو۔ نو $!C_1$ به بيا منځګړی وي، که دا شرط پوره
    وي چې $\Lang L(!C_1) \subseteq \Lang L(\Gamma_1, \Delta_1,
    \lexists[x][!A(x)]) \cap \Lang L(\Gamma_2, \Delta_2)$۔ دلته
    استقرايي فرض $\Lang L(!C_1)
    \subseteq \Lang L(\Gamma_1, \Delta_1, \lexists[x][!A(x)], !A(t))
    \cap \Lang L(\Gamma_2, \Delta_2)$ راکوي۔ اوس ياد کړئ چې
    فرض مو کړے ؤ چې هېڅ !!{function}s نۀ شته۔ نو د دې قواعدو
    تړلے ترم~$t$ يوازې !!a{constant} کېدے شي؛ يعنې $t \ident b$۔
    که $b \notin \Lang L(!C_1)$ يا $b \in \Lang L(\Gamma_1, \Delta_1, \lexists[x][!A(x)]) \cap
    \Lang L(\Gamma_2, \Delta_2)$ وي، اندېښنه نۀ شته:
    $!C_1$ منځګړی اخستلے شو۔ خو که $b \in \Lang
    L(!C_1)$ او $b \notin \Lang L(\Gamma_1, \Delta_1,
    \lexists[x][!A(x)]) \cap \Lang L(\Gamma_2, \Delta_2)$ وي، بيا؟
    لومړے، په~$!C_1$ کښې د~$b$ هر راتګ په يوه تازه
    !!a{variable}~$y$ بدلولے شو او $!C_1 \ident \Subst{!C_1'}{b}{y}$
    ليکلے شو، چې~$!C_1'$ پکښې~$b$ نۀ لري۔ سربېره پردې، يا
    $b \notin
    \Lang L(\Gamma_1, \Delta_1, \lexists[x][!A(x)])$ يا $b \notin
    \Lang L(\Gamma_2, \Delta_2)$ دے۔ په ترتيب سره $!C \ident
    \lforall[y][!C_1'(y)]$ يا $!C \ident \lexists[y][!C_1'(y)]$
    اخستلے شو۔ په لومړي حالت کښې لرو:
\[
    \AxiomC{}
    \Deduce$\Gamma_1 \fCenter
      \Delta_1, \lexists[x][!A(x)], !A(b), !C_1'(b)$
    \RightLabel{\RightR{\lexists}}
    \UnaryInf$\Gamma_1 \fCenter \Delta_1, \lexists[x][!A(x)], !C_1'(b)$
    \RightLabel{\RightR{\lforall}}
    \UnaryInf$\Gamma_1 \fCenter
      \Delta_1, \lexists[x][!A(x)], \lforall[y][!C_1'(y)]$
    \DisplayProof
    \]
    او
\[
    \AxiomC{}
    \Deduce$!C_1'(b), \lforall[y][!C_1'(y)], \Gamma_2 \fCenter \Delta_2$
    \RightLabel{\LeftR{\lforall}}
    \UnaryInf$\lforall[y][!C_1'(y)], \Gamma_2 \fCenter \Delta_2$
    \DisplayProof
    \]
    د پاسنيو سېکوېنټونو !!{proof}s له استقرايي فرضه راځي؛ په
    دويم کښې په مقدم کښې د $\lforall[y][!C_1'(y)]$ د زياتولو
    دپاره د کمزوري کولو منل کېدل کاروو۔ په لومړي !!{proof}
    کښې د \RightR{\lforall} د ځانګړي متغير شرط پوره دے، ځکه
    $b \notin \Lang L(\Gamma_1, \Delta_1,
    \lexists[x][!A(x)])$ او~$!C_1'$ داسې تعريف شوے چې~$b$
    نۀ لري۔

    په دويم حالت کښې، له استقرايي فرض او د~$\lexists[y][!C_1'(y)]$
    په زياتولو د کمزوري کولو له منل کېدو څخه، لرو:
\[
    \AxiomC{}
    \Deduce$\Gamma_1 \fCenter
      \Delta_1, \lexists[x][!A(x)], !A(b), \lexists[y][!C_1'(y)], !C_1'(b)$
    \RightLabel{\RightR{\lexists}}
    \UnaryInf$\Gamma_1 \fCenter
      \Delta_1, \lexists[x][!A(x)], \lexists[y][!C_1'(y)], !C_1'(b)$
    \RightLabel{\RightR{\lexists}}
    \UnaryInf$\Gamma_1 \fCenter
      \Delta_1, \lexists[x][!A(x)], \lexists[y][!C_1'(y)]$
    \DisplayProof
    \]
    او
\[
    \AxiomC{}
    \Deduce$!C_1'(b), \Gamma_2 \fCenter \Delta_2$
    \RightLabel{\LeftR{\lexists}}
    \UnaryInf$\lexists[y][!C_1'(y)], \Gamma_2 \fCenter \Delta_2$
    \DisplayProof
    \]
    په دويم !!{proof} کښې د \LeftR{\lexists} د ځانګړي متغير
    شرط پوره دے، ځکه $b \notin \Lang L(\Gamma_2,
    \Delta_2)$ او $!C_1'$ داسې تعريف شوے چې~$b$ نۀ لري۔
  \end{enumerate}
  نور حالتونه هم ورته دي؛ تمرينونو ته ئې پرېږدو۔
\end{proof}

\begin{prob}
فرض کړئ چې د $\lnand$ په نوم يو تړونکے د دې قواعدو سره لرو:
\[
\Axiom$\Gamma \fCenter \Delta, !A$
\Axiom$\Gamma \fCenter \Delta, !B$
\RightLabel{\LeftR{\lnand}}
\BinaryInf$!A \lnand !B, \Gamma \fCenter \Delta$
\DisplayProof
\quad
\Axiom$!A, !B, \Gamma \fCenter \Delta$
\RightLabel{\RightR{\lnand}}
\UnaryInf$\Gamma \fCenter \Delta, !A \lnand !B$
\DisplayProof
\]
د \olref{lem:maehara} د ثبوت هغه حالتونه ترسره کړئ چې $\pi$
په دغو قواعدو کښې په يوه ختمېږي، او وښيئ چې د مايېهارا
مرستندويه قضيه لا هم صدق کوي۔
\end{prob}

اوس چې د مايېهارا مرستندويه قضيه مو ثابته کړه، د کرېګ د
منځګړيتوب د قضيې د ثابتولو دپاره ئې کارولے شو۔

\begin{proof}[د \olref{thm:interpolation} ثبوت]
  فرض کړئ چې $!A \Sequent !B$ !!{provable} دے۔ په
  $\maeh{!A}{\ } \Sequent \maeh{\ }{!B}$ باندې \olref{lem:maehara}
  وکاروئ، څو منځګړی~$!C$ تر لاسه کړئ۔ $\Proves !A \Sequent !C$
  او $\Proves !C
  \Sequent !B$ لرو۔
\end{proof}

فرض کړئ چې $\Gamma$ په !!a{language}~$\Lang L$ کښې يوه تيوري
(د !!{sentence}s سټ) ده او ژبه د $n$ ځايونو {مسند}~$R$ لري۔
وايو چې $\Gamma$،~$R$ \emph{په ضمني توګه تعريفوي}، هله او يوازې هله چې
\begin{align*}
\Gamma, \Gamma'
& \Entails \lforall[x_1][\dots\lforall[x_n][(R(x_1,\dots,x_n) \liff
R'(x_1, \dots, x_n))]],
\intertext{دلته $\Gamma'$ له $\Gamma$ څخه د~$R$ هر راتګ په~$R' \notin \Lang L$
بدلولو تر لاسه کېږي۔ $\Gamma$،~$R$ \emph{په څرګند ډول تعريفوي}،
هله او يوازې هله چې داسې~$!A(x_1, \dots, x_n)$ شته چې~$R$ نۀ لري او}
\Gamma &\Entails
\lforall[x_1][\dots\lforall[x_n][(R(x_1,\dots,x_n) \liff !A(x_1,
\dots, x_n))]].
\end{align*}
څرګنده ده: که $\Gamma$، $R$ په څرګند ډول تعريفوي، نو~$R$ په
ضمني توګه هم تعريفوي۔ معکوسه خبره څرګنده نۀ ده، خو بيا هم صدق کوي:

\begin{lem}[د بېت د تعريف‌وړتيا قضيه]
  که $\Gamma$،~$R$ په ضمني توګه تعريفوي، نو په څرګند ډول ئې هم
  تعريفوي۔
\end{lem}

\begin{proof}
  فرض کړئ چې $\Gamma$،~$R$ په ضمني توګه تعريفوي۔ دې $R'$ او
  $\Gamma'$ د پورته په شان وي او $\Lang L' = \Lang L(\Gamma') = \Lang L \setminus \{R\}
  \cup \{R'\}$ دې وي۔ د فرض له مخې لرو:
\begin{align*}
    \Gamma, \Gamma' & 
    \Sequent \lforall[x_1][\dots\lforall[x_n][(R(x_1,\dots,x_n) 
    \liff R'(x_1, \dots, x_n))]]
  \intertext{دې $c_1$، \dots،~$c_n$ داسې !!{constant}s وي چې په~$\Gamma$ کښې نۀ راځي۔ د معکوسولو له مرستندويې قضيې څخه لرو}
    \Gamma, \Gamma', R(c_1,\dots,c_n) & \Sequent  R'(c_1, \dots, c_n)
  \intertext{د مايېهارا مرستندويه قضيه په دې وکاروئ}
    \maeh{\Gamma, R(c_1,\dots,c_n)}{\Gamma'} & \Sequent
      \maeh{\ }{R'(c_1, \dots, c_n)}
  \intertext{څو داسې~$!A \ident !A(x_1, \dots, x_n)$ ومومو چې}
    \Gamma, R(c_1,\dots,c_n) & \Sequent !A(c_1, \dots, c_n) \\
    !A(c_1, \dots, c_n), \Gamma' & \Sequent R'(c_1, \dots, c_n)
  \intertext{!!{proof}s لري۔ \textbf{سرچينه‌يي سمون OLCMP-121:} مرستندويه قضيه د ګډې ژبې د غړيتوب حد راکوي، نۀ لازمي برابري۔ موقتي منځګړی تازه ثابتونه لرلے شي؛ د هغو په ځاے متغيرونو اخستلو وروسته $\Lang L(!A) \subseteq \Lang L(\Gamma) \setminus \{R\}$ دے، نو~$!A$،~$R$ نۀ لري۔ په دويم سېکوېنټ کښې د !!{proof} په هر ځاے $R'$ په~$R$ بدل کړئ، څو د لاندې سېکوېنټ !!a{proof} تر لاسه شي}
    !A(c_1, \dots, c_n), \Gamma & \Sequent R(c_1, \dots, c_n)
  \end{align*}
  \RightR{\lif} او \RightR{\lforall} وکاروئ، څو د لاندې
  سېکوېنټ !!a{proof} تر لاسه کړئ:
\[
    \Gamma \Sequent \lforall[x_1][\dots\lforall[x_n][(R(x_1,\dots,x_n) \liff !A(x_1, \dots, x_n))]].
  \]
  دا ښيي چې $!A(x_1, \dots, x_n)$ په~$\Gamma$ کښې د $!R$
  څرګند تعريف دے۔
\end{proof}

درېيمه پايله، چې منځګړيتوب ته هم‌ارزښته ده او د هغې په شان
د مايېهارا په مرستندويه قضيه ثابتېدے شي، دا قضيه ده: که دوه
سازګارې تيورۍ~$\Gamma_1$ او $\Gamma_2$ د خپلې ټولې ګډې ژبې
يوه بشپړه تيوري $\Gamma$ ولري، يعنې
$\Lang L(\Gamma) = \Lang L(\Gamma_1)
\cap \Lang L(\Gamma_2)$ وي، نو $\Gamma_1 \cup \Gamma_2$ سازګاره ده۔
\textbf{سرچينه‌يي سمون OLCMP-122:} يوازې د ګډې ژبې په يوه
کوچنۍ فرعي ژبه کښې بشپړتيا کافي نۀ ده؛ لاندې استدلال د ټولې
ګډې ژبې د هرې جملې فيصله غواړي۔

په ياد ولرئ چې بشپړه تيوري هغه ده چې په خپله ژبه کښې د هرې
!!{sentence}~$!A$ دپاره يا $\Gamma \Proves !A$ يا
$\Gamma \Proves \lnot !A$ وي۔ ځکه چې $\Gamma_1$ او~$\Gamma_2$
د~$\Gamma$ سازګارې پراختياوې دي، نو $\Gamma$ بايد سازګاره وي۔
له دې راځي: که $!A$ د $\Lang L(\Gamma_1) \cap \Lang
L(\Gamma_2)$ په ژبه کښې !!a{sentence} وي، نو $\Gamma_1 \Entails !A$
او $\Gamma_2 \Entails \lnot !A$ دواړه نۀ شي لرلے۔ د دې د کتلو
دپاره فرض کړئ داسې~$!A$ شته۔ که $\Gamma_1 \Entails !A$ وي، نو
$\Gamma \Entails !A$ دے۔ که داسې نۀ وي، د~$\Gamma$ د بشپړتيا
له مخې $\Gamma \Entails \lnot !A$ دے؛ او ځکه چې
$\Gamma \subseteq \Gamma_1$ ده، $\Gamma_1 \Entails \lnot !A$
بۀ وي او $\Gamma_1$ به ناسازګاره شي۔ نو $\Gamma \Entails !A$
دے؛ او ځکه چې $\Gamma
\subseteq \Gamma_2$ ده، $\Gamma_2 \Entails !A$ دے۔ نو
$\Gamma \Entails/
\lnot !A$، ګنې $\Gamma_2$ به ناسازګاره وي۔

د ثبوت دپاره همدا کافي ده چې په ګډه !!{language}~$\Lang L(\Gamma_1) \cap \Lang L(\Gamma_2)$
کښې داسې !!a{sentence}~$!A$ نۀ وي چې $\Gamma_1
\Entails !A$ او $\Gamma_2 \Entails \lnot !A$ وي۔ دلته دا
ثبوت د ثبوت پوهنې په طريقه، د مايېهارا په مرستندويه قضيه، ورکوو۔

\begin{thm}[د رابنسن د ګډې سازګارۍ قضيه]
فرض کړئ چې $\Gamma_1$ او $\Gamma_2$ سازګارې تيورۍ دي او په
!!{language} $\Lang L(\Gamma_1) \cap \Lang L(\Gamma_2)$ کښې
د هېڅ~$!A$ دپاره $\Gamma_1 \Entails !A$ او
$\Gamma_2 \Entails \lnot !A$ دواړه نۀ لرو۔ نو
$\Gamma_1 \cup \Gamma_2$ سازګاره ده۔
\end{thm}

\begin{proof}
  د خلاف فرض وکړئ چې $\Gamma_1 \cup \Gamma_2$ ناسازګاره ده۔
  نو سېکوېنټ $\Gamma_1, \Gamma_2 \Sequent \ $ يو
  !!a{proof} لري۔ په $\maeh{\Gamma_1}{\Gamma_2}
  \Sequent \maeh{\ }{\ }$ د مايېهارا مرستندويه قضيه وکاروئ،
  څو په !!{language}~$\Lang L(\Gamma_1) \cap \Lang L(\Gamma_2)$
  کښې داسې !!a{sentence}~$!A$ تر لاسه کړئ چې
  $\Proves \Gamma_1 \Sequent !A$ او $\Proves !A, \Gamma_2 \Sequent
  \quad$ وي۔ له وروستي څخه $\Proves \Gamma_2 \Sequent \lnot !A$
  تر لاسه کوو۔ خو فرض مو کړے ؤ چې داسې !!{sentence} نۀ شته۔
\end{proof}

پورته مو په پټه فرض کړے ؤ چې تيورۍ $\Gamma$، $\Gamma_1$ او
$\Gamma_2$ متناهي دي۔ د متناهي شاهد خاصيت له مخې، پايلې د
نامتناهي تيوريو دپاره هم صدق کوي۔

\end{document}
